% =====================================================================
%  K-LD :: E3 -- Differential validation of the ESBMC LD->GOTO translation
%  Self-contained; uses the template idiom (\gls, \code, booktabs, \paragraph).
% =====================================================================
\providecommand{\kld}{\textsc{K-LD}}
\providecommand{\openplc}{OpenPLC}
\providecommand{\cmark}{\ensuremath{\checkmark}}      % amssymb
\providecommand{\xmark}{\ensuremath{\boldsymbol{\times}}}

\section{Differential Validation of the LD\texorpdfstring{$\rightarrow$}{->}GOTO Translation}
\label{sec:e3}

Having established that \kld{} is a faithful executable model of \gls{iec}~61131-3
(\S\ref{sec:eval:rq1}), we use it as a \emph{reference oracle} for the \gls{esbmc}
\gls{ld}$\rightarrow$GOTO translation. This section answers two questions: on the
benchmark suite, does \gls{esbmc} agree with \kld{} (RQ2), and where they disagree,
which engine is correct and why (RQ3).

\subsection{Experimental Setup}
\label{sec:e3:setup}

\paragraph{Three engines.} Each benchmark is checked by (i)~\gls{esbmc}~v8.3.0 via its
\gls{ld} front-end (\code{--ld-props}, \code{--unwind 6 --no-unwinding-assertions});
(ii)~the \kld{} interpreter; and, for any disagreement, (iii)~\openplc{}'s runtime
(\gls{matiec}-generated C), whose reference function-block semantics \kld{} was
validated against scan-for-scan in \S\ref{sec:eval:rq1}. Engine~(iii) is the external
tie-breaker: it decides which of \gls{esbmc} and \kld{} matches real \gls{plc}
execution.

\paragraph{Two independent front-ends.} To keep \kld{} an \emph{independent} oracle we
never reuse \gls{esbmc}'s parser. A translator lowers each PLCopen~XML benchmark into
\kld{} \gls{dsl}: a \emph{simple} path for the linear \code{<rung>} format (contacts in
series~$=$~AND, coils' \code{storage} set/reset~$\rightarrow$~latch/unlatch) and a
\emph{graphical} path that resolves the netlist (power rails, series/parallel wiring,
\code{TON}/\code{TOF}/\code{TP} and \code{CTU}/\code{CTD} blocks, and rising/falling
edge contacts lowered to \code{R\_TRIG}/\code{F\_TRIG} helper logic). The translators
carry no verification semantics; they only re-express the diagram.

\paragraph{Bounded differential.} Because \gls{esbmc} verifies over a bounded scan
horizon, \kld{} is exercised comparably: we drive every input combination (a bounded
random sample when the input space is large) across several scan-cycle sweeps and
evaluate each property on every scan of the resulting execution. Properties are
interpreted exactly as \gls{esbmc}'s \code{--ld-props} does, honouring the
\code{kind} field: an \code{invariant} must hold in every reachable image, an
\code{absence} condition must never hold, and \code{mutual\_exclusion} of a variable
set forbids any two being true at once. A program is \textsc{unsafe} for \kld{} iff
some scan falsifies some property, with that scan reported as a witness.

\subsection{RQ2: \kld{} Agrees with \texorpdfstring{\gls{esbmc}}{ESBMC} on Well-Modelled Programs}
\label{sec:e3:rq2}

Table~\ref{tab:e3} reports all 13 benchmark programs. On \textbf{10 of 13} the two
engines return the same verdict. The agreements span the full construct fragment:
combinational interlocks (\code{tank\_level\_control}), sealed-in logic
(\code{water\_control}), and timer/edge-driven controllers (\code{bottle\_filling},
\code{elevator}, \code{beremiz\_traffic\_light}). Crucially, on every \textsc{unsafe}
variant \kld{} not only matches \gls{esbmc}'s pass/fail but \emph{localises} the
violation to specific properties with a concrete witness image---for example, the
unsafe \code{bottle\_filling} flags \code{P1}\,(valve open with no bottle) and
\code{P2}\,(valve open when full), while correctly keeping \code{P3},\code{P5} safe
because the emergency-stop guard still propagates through \code{System\_Running} within
a scan. This per-property, witness-bearing verdict is strictly more informative than
\gls{esbmc}'s aggregate result.

\begin{table}[t]
  \centering
  \caption{Three-engine differential over the benchmark suite. \kld{} verdicts list
    the violated properties. Rows marked \textbf{\xmark} disagree; all three are
    resolved in \kld{}'s favour in \S\ref{sec:e3:rq3}.}
  \label{tab:e3}
  \small
  \begin{tabular}{@{}llll c@{}}
    \toprule
    Program & Fmt & \kld{} & \gls{esbmc} & Agree \\
    \midrule
    tank\_level\_control          & rung & \textsc{safe}                 & \textsc{safe}   & \cmark \\
    tank\_level\_control\,(unsafe)& rung & \textsc{unsafe} P1,P2,P3      & \textsc{unsafe} & \cmark \\
    bottle\_filling               & rung & \textsc{safe}                 & \textsc{safe}   & \cmark \\
    bottle\_filling\,(unsafe)     & rung & \textsc{unsafe} P1,P2,P4,P6   & \textsc{unsafe} & \cmark \\
    elevator                      & rung & \textsc{safe}                 & \textsc{safe}   & \cmark \\
    elevator\,(unsafe)            & rung & \textsc{unsafe} P1--P6        & \textsc{unsafe} & \cmark \\
    traffic\_light\,(safe)        & rung & \textsc{safe}                 & \textsc{safe}   & \cmark \\
    water\_control                & graph & \textsc{safe}                & \textsc{safe}   & \cmark \\
    beremiz\_traffic\_light       & graph & \textsc{safe}                & \textsc{safe}   & \cmark \\
    dimmer\_light\_control        & graph & \textsc{safe}                & \textsc{safe}   & \cmark \\
    \midrule
    \textbf{stairs\_light}        & graph & \textbf{\textsc{unsafe} P1}   & \textsc{safe}   & \xmark \\
    \textbf{traffic\_light}       & rung & \textsc{safe}                 & \textbf{\textsc{unsafe}} & \xmark \\
    \textbf{traffic\_light\,(unsafe)} & rung & \textsc{safe}             & \textbf{\textsc{unsafe}} & \xmark \\
    \bottomrule
  \end{tabular}
\end{table}

\subsection{RQ3: Every Disagreement Is an \texorpdfstring{\gls{esbmc}}{ESBMC} Timer Defect}
\label{sec:e3:rq3}

The three disagreements all involve timer function blocks, and \openplc{} confirms
\kld{} in every case. They expose two \emph{opposite} defects in \gls{esbmc}'s
incomplete timer translation (Table~\ref{tab:e3tax}).

\paragraph{A controlled probe.} We first isolate the mechanism with a minimal program:
\code{Btn}$\rightarrow$\code{TON}$(PT)$$\rightarrow$\code{Light}. Under a faithful timer
the invariant \code{A}: ``\code{!Light || Btn}'' holds---\code{Light} follows the done
bit, and a \code{TON} holds its output only while its input is high. \kld{} proves
\code{A} safe; \gls{esbmc} reports it \emph{violated}, with a counterexample in which
the done bit is true while the input is false---a state a \code{TON} can never reach,
since a low input resets it. \gls{esbmc} therefore leaves the timer output
\emph{unconstrained} (havoc), an over-approximation that manufactures false alarms.

\paragraph{Discrepancy 1---havoc false alarms (\code{traffic\_light}).} The real
\code{traffic\_light} controller sequences four phases through a chain of four
\code{TON}s. \gls{esbmc} declares it \textsc{unsafe} on the green-conflict
mutual-exclusion property; its counterexample asserts \emph{all four phases active in
the same scan}. A real \code{TON} chain fires one phase at a time, so this state is
unreachable---it arises only because \gls{esbmc} havocs the four done bits. Executed
under \kld{}, the controller exhibits \emph{zero} mutual-exclusion violations across a
full multi-cycle horizon (indeed it deadlocks in the north--south phase, a liveness
defect that violates no \emph{safety} property). \kld{}'s \textsc{safe} verdict is
correct, and independent of the timer preset.

\paragraph{Discrepancy 2---skipped timers miss real bugs (\code{stairs\_light}).} The
graphical \code{stairs\_light} benchmark drives its light through an off-delay
(\code{TOF}) so that a single PIR detection keeps the light on for the preset. Its
front-end warns that ``FB/timer outputs on rail$\rightarrow$coil paths are not yet
modelled'' and \emph{skips} the timer path, after which \gls{esbmc} certifies the
program \textsc{safe}. But that skip is exactly where the violation lives: once the PIR
clears while the light is still held on by the off-delay, invariant \code{P1}
(``light on implies PIR active or button state on'') is falsified. \kld{}, modelling
the \code{TOF}, reports \code{P1} \textsc{unsafe} with the light-on/PIR-off witness, and
\openplc{} reproduces the same light behaviour. Here \gls{esbmc} is \emph{unsound}: it
returns \textsc{safe} for a program that violates its own stated property.

\begin{table}[t]
  \centering
  \caption{Discrepancy taxonomy. Every disagreement is an \gls{esbmc} timer-translation
    defect; \kld{} is correct in all three, corroborated by \openplc{}.}
  \label{tab:e3tax}
  \small
  \begin{tabular}{@{}lllp{3.5cm}@{}}
    \toprule
    Program & \gls{esbmc} & \kld{} & \gls{esbmc} defect \\
    \midrule
    stairs\_light   & \textsc{safe}   & \textsc{unsafe} & \emph{skips} the FB path (graphical) $\rightarrow$ under-approx.\ $\rightarrow$ \textbf{missed bug} (unsound) \\
    traffic\_light  & \textsc{unsafe} & \textsc{safe}   & \emph{havocs} the FB output (rung) $\rightarrow$ over-approx.\ $\rightarrow$ \textbf{false alarm} \\
    \bottomrule
  \end{tabular}
\end{table}

\subsection{Discussion and Threats}
\label{sec:e3:threats}

\paragraph{What the oracle bought.} The differential is decisive precisely because
\kld{} is an \emph{executable} \gls{iec}-faithful semantics: it does not merely flag a
mismatch, it produces the concrete scan-level behaviour that shows which engine is
right, and \openplc{} confirms it. The two defects are complementary and both
practically serious---the havoc mode wastes engineering effort on impossible
counterexamples, while the skip mode is unsound and can certify a hazardous program safe.

\paragraph{Threats to validity.} (i)~The comparison is bounded; \kld{}'s \textsc{safe}
results are safe-up-to-horizon, matching \gls{esbmc}'s own bounded model, and the
\code{traffic\_light} refutation of \gls{esbmc} is preset-independent by construction.
(ii)~Timer presets absent from a benchmark are instantiated with a fixed default; this
affects \emph{when} a timer fires, not the existence of the discrepancies, each of
which we tie to a structural argument. (iii)~The two front-ends are independent of
\gls{esbmc} and were cross-checked by reproducing a manually translated benchmark
identically, so a divergence cannot be an artefact of a shared parser.
